\chapter{Introducción}

\section{Contexto}

La propuesta MuseIQ se inscribe en la convergencia entre transformación digital, mediación cultural y accesibilidad en instituciones museísticas. En el Perú, la adopción tecnológica en hogares configura un escenario favorable para soluciones móviles de apoyo a la visita. La tenencia de smartphones evolucionó de 78.0\% en 2019 a 94.8\% en 2024, y en 2023 se registraron 9\,779\,068 hogares con este dispositivo~\cite{osiptel2023smartphone}. Para 2024, OSIPTEL reportó que la cifra superó los 10 millones de familias~\cite{osiptel2024smartphone}. Esta trayectoria de adopción confirma la viabilidad de usar el teléfono móvil como interfaz principal para servicios de guía inteligente basados en proximidad y contexto.

\vspace{0.5cm}

\begin{table}[h!]
	\centering
	\begin{tabular}{l c c}
		\hline
		\textbf{Año} & \textbf{Indicador oficial} & \textbf{Dato reportado} \\
		\hline
		2019 & Hogares con smartphone & 78.0 \% \\
		2021 & Hogares con smartphone & 88.4 \% \\
		2022 & Hogares con smartphone & 91.9 \% \\
		2023 & Hogares con smartphone & 92.8 \% \\
		2024 & Hogares con smartphone & 94.8 \% \\
		\hline
	\end{tabular}
	\caption{Evolución de hogares peruanos con smartphone \textbf{Fuente:}~\cite{osiptel2021smartphone,osiptel2022smartphone,osiptel2023smartphone,osiptel2024smartphone}}
	\label{tab:smartphones-peru-2019-2024}
\end{table}

En paralelo, el sector museístico peruano evidencia una recuperación sostenida de la demanda presencial. El Plan Estratégico Sectorial Multianual 2022--2030 reportó 1\,956\,034 visitas en 2019~\cite{culturapesem2021}; durante 2020, la restricción sanitaria redujo la cifra a 353\,153~\cite{cultura2020memoria}. Posteriormente, la recuperación alcanzó 1\,067\,571 visitas en 2022~\cite{cultura2022memoria}, 1\,543\,872 en 2023~\cite{cultura2023memoria} y 1\,790\,502 en 2024~\cite{cultura2024reporte}. Aunque el nivel de 2024 todavía se ubica por debajo del valor prepandemia, la tendencia ascendente confirma masa crítica de usuarios para la implementación y validación de nuevas soluciones de mediación cultural.

\vspace{0.5cm}

\begin{table}[h!]
	\centering
	\begin{tabular}{l c l}
		\hline
		\textbf{Año} & \textbf{Visitantes} & \textbf{Observación} \\
		\hline
		2019 & 1\,956\,034 & Línea base prepandemia \\
		2020 & 353\,153 & Contracción por COVID-19 \\
		2022 & 1\,067\,571 & Recuperación parcial \\
		2023 & 1\,543\,872 & Recuperación sostenida \\
		2024 & 1\,790\,502 & Máximo pospandemia, aún por debajo de 2019 \\
		\hline
	\end{tabular}
	\caption{Evolución de visitas a museos administrados por el Ministerio de Cultura \textbf{Fuente:}~\cite{culturapesem2021,cultura2020memoria,cultura2022memoria,cultura2023memoria,cultura2024reporte}}
	\label{tab:visitas-museos-peru-2019-2024}
\end{table}

En el plano internacional, UNESCO plantea que las tecnologías digitales están redefiniendo las condiciones de acceso, experiencia y preservación de bienes culturales, y enfatiza su potencial para reducir desigualdades y promover innovación sectorial~\cite{unesco2025globalreport}. Esta evidencia respalda la pertinencia de MuseIQ como respuesta aplicada a una transformación estructural de la gestión cultural contemporánea.

\section{Descripción del Problema}

Pese al crecimiento de la demanda y a la disponibilidad de dispositivos móviles, una proporción relevante de museos mantiene mecanismos de mediación predominantemente estáticos, tales como cartelas físicas, recorridos lineales y audioguías poco adaptativas. Esta configuración limita la entrega de información situada, reduce la capacidad de personalización según perfil o ritmo de visita y mantiene brechas de accesibilidad para usuarios con requerimientos específicos~\cite{ivanov2025smartmuseums,wang2024voiceguide}.

Desde una perspectiva operativa, el problema central se formula del siguiente modo: los museos requieren una solución de mediación digital capaz de estimar con confiabilidad suficiente la ubicación y orientación del visitante dentro de la sala, activar contenido contextual y sostener costos de implementación compatibles con restricciones presupuestales institucionales. La persistencia de esta brecha entre demanda cultural creciente y digitalización parcial reduce el potencial educativo, inclusivo y de apropiación social de la visita presencial.

\section{Solución Tecnológica}

La solución propuesta consiste en una guía inteligente para museos articulada sobre cuatro componentes: beacons BLE basados en ESP32, sensores del smartphone, interacción por voz mediante TTS/STT y una capa de recuperación aumentada para generación de respuestas (RAG). El objetivo técnico no es el posicionamiento centimétrico, sino la identificación robusta de sala y zona probable de observación para activar contenido pertinente en tiempo real.

El uso de BLE se justifica por su relación costo-beneficio y su escalabilidad operativa. Un caso implementado en el Foz C\^oa Museum reportó 96\% de precisión en identificación de sala mediante BLE y RSSI~\cite{verde2023blemuseum}. Esta evidencia valida la factibilidad de un esquema de reconocimiento por ambiente en contextos museísticos que no disponen de infraestructura de alto costo.

Asimismo, la documentación oficial de Android establece que la orientación del dispositivo puede inferirse combinando acelerómetro y sensor geomagnético~\cite{android2026positionsensors}. Esta información complementa la señal de proximidad BLE y fortalece la activación contextual de contenidos. Finalmente, la capa RAG permite consultar bases curatoriales institucionales antes de generar respuestas, reduciendo alucinaciones y aumentando pertinencia semántica~\cite{aws2026rag}.

\section{Referencias a tomar en cuenta}

\subsection{Indoor Content Delivery Solution for a Museum Based on BLE Beacons}

El estudio de Verde et~al. propone y valida una arquitectura de posicionamiento y entrega de contenido en interiores mediante BLE en un entorno museístico real~\cite{verde2023blemuseum}. Su aporte principal para esta investigación radica en demostrar, con evidencia empírica, que la identificación de sala puede alcanzar alta precisión sin infraestructura compleja.

\subsection{Co-design of a Voice-Driven Interactive Smart Guide for Museum Accessibility and Management}

La investigación de Wang enfatiza la interacción por voz en guías inteligentes como mecanismo de accesibilidad y mejora de experiencia de usuario~\cite{wang2024voiceguide}. Este antecedente sustenta la incorporación de TTS/STT en MuseIQ dentro de una estrategia de inclusión funcional en recorrido presencial.

\subsection{Analyzing Visitor Behavior to Enhance Personalized Experiences in Smart Museums}

La revisión sistemática de Ivanov y Velkova sintetiza evidencia sobre personalización en museos inteligentes y destaca el papel del seguimiento de ubicación, la analítica de comportamiento y la adaptación en tiempo real~\cite{ivanov2025smartmuseums}. Este trabajo provee sustento conceptual para diseñar experiencias no lineales centradas en el visitante.

\subsection{Referentes institucionales y de mercado}

Las series oficiales de OSIPTEL y del Ministerio de Cultura permiten dimensionar, con base empírica, el contexto nacional de adopción tecnológica y demanda de visitas~\cite{osiptel2021smartphone,osiptel2022smartphone,osiptel2023smartphone,osiptel2024smartphone,culturapesem2021,cultura2020memoria,cultura2022memoria,cultura2023memoria,cultura2024reporte}. Adicionalmente, Bloomberg Connects evidencia la consolidación de soluciones móviles de mediación cultural a escala internacional~\cite{bloomberg2025connects}.

\vspace{0.5cm}

\begin{table}[h!]
	\centering
	\small
	\setlength{\tabcolsep}{3.5pt}
	\begin{tabular}{>{\raggedright\arraybackslash}p{4.3cm} >{\raggedright\arraybackslash}p{2.5cm} >{\raggedright\arraybackslash}p{3.6cm} >{\raggedright\arraybackslash}p{3.6cm}}
		\hline
		\textbf{Referente} & \textbf{Tipo} & \textbf{Aporte principal} & \textbf{Relevancia para MuseIQ} \\
		\hline
		Verde et al. (2023), \textit{Indoor Content Delivery Solution for a Museum Based on BLE Beacons} & Artículo científico & Implementa BLE + RSSI en museo real; reporta 96\% de precisión de sala & Valida reconocimiento por sala y activación contextual \\
		Wang (2024), \textit{Co-design of a Voice-Driven Interactive Smart Guide for Museum Accessibility and Management} & Artículo científico & Diseña guía por voz centrada en accesibilidad y experiencia de visitante & Sustenta uso de TTS/STT y enfoque inclusivo \\
		Ivanov y Velkova (2025), \textit{Analyzing Visitor Behavior to Enhance Personalized Experiences in Smart Museums} & Revisión sistemática & Sintetiza evidencia sobre personalización en tiempo real y tracking & Sustenta la lógica de personalización contextual \\
		Bloomberg Connects & Referente de mercado & Escalamiento internacional de guías móviles en instituciones culturales & Evidencia demanda y madurez del mercado objetivo \\
		\hline
	\end{tabular}
	\caption{Trabajos y referentes clave para el desarrollo de MuseIQ \textbf{Fuente:}~\cite{verde2023blemuseum,wang2024voiceguide,ivanov2025smartmuseums,bloomberg2025connects}}
\end{table}

\section{Estado del Arte}

El estado del arte en museos inteligentes converge en tres líneas de desarrollo: localización en interiores, personalización basada en comportamiento y experiencias multimodales accesibles. La literatura reciente indica que estas líneas alcanzan mayor efectividad cuando se diseñan como una arquitectura integrada orientada al desempeño de la experiencia de visita~\cite{ivanov2025smartmuseums}.

En localización indoor, BLE destaca como alternativa pragmática frente a tecnologías de mayor costo. La evidencia en museos reales muestra que la identificación correcta de sala aporta valor operacional suficiente para activar contenido contextual, incluso sin perseguir precisión milimétrica~\cite{verde2023blemuseum}. Esta conclusión resulta crítica para un enfoque MVP, donde costo, mantenimiento y escalabilidad constituyen restricciones de diseño.

En accesibilidad, la interacción por voz amplía la cobertura funcional para usuarios con barreras visuales o de lectura sostenida, y mejora la fluidez de uso durante recorridos presenciales~\cite{wang2024voiceguide}. En paralelo, los ecosistemas de guías móviles a escala internacional confirman una demanda sostenida por experiencias culturales digitales en dispositivos personales~\cite{bloomberg2025connects}.

En consecuencia, el aporte diferencial de esta tesis consiste en proponer una arquitectura unificada, de bajo costo y contextualizada al entorno peruano, que integra reconocimiento por sala mediante BLE, apoyo de orientación con sensores móviles e interacción conversacional sustentada en recuperación de conocimiento institucional. La contribución no se limita a un componente aislado, sino a la articulación técnica de dichos elementos para resolver una brecha concreta de mediación cultural presencial.
